\documentclass[12pt,aspectratio=169,ignorenonframetext]{ltxbeamer}

\usepackage[spanish,es-nodecimaldot,es-tabla]{babel}
\usepackage{tabularx}
\usepackage{multirow}

\title{Estructura}
\subtitle{Curso de {\lmr\LaTeX}}
\author{Joar Buitrago\texorpdfstring{\thanks{jebuitragoc@unal.edu.co}}{ <jebuitragoc@unal.edu.co>}}
\date{}

\sisetup{
    unit-mode=text,
    text-font-command=\lmr,
}

\addbibresource{bibliography.bib}



\begin{document}
\begin{frame}
\maketitle
\end{frame}

\section{Estructura}

\begin{frame}[fragile]
\frametitle{Comandos}

Todos los comandos en {\lmr\LaTeX} se escriben con una barra inversa (\texttt{\textbackslash}) seguidos de caracteres alfabéticos.

\pause

Todos los comandos pueden tener argumentos, los cuales son distinguidos entre opcionales (\mintinline{latex}{[...]}) y obligatorio (\mintinline{latex}{{...}}).

\pause

\begin{exampleblock}{}
    \begin{minted}[escapeinside=||]{latex}
        \item
        \raisebox{<|\textit{distance}|>}[<|\textit{height}|>][<|\textit{depth}|>]{<|\textit{text}|>}
    \end{minted}
\end{exampleblock}
\end{frame}





\begin{frame}[fragile]
\frametitle{Tipos de Comandos}

Existen varios tipos de comandos definidos en {\lmr\LaTeX}:

\pause

\begin{description}[<+->]
    \item[Simples] Son comandos sin argumentos que se usan generalmente para producir símbolos o efectos.
    \item[Básicos] Son comandos con uno o más argumentos.
    \item[Globales] Son comandos, generalmente simples, que afectan el comportamiento de {\lmr\LaTeX}. Este comportamiento permanece vigente dentro del ámbito en el que se utilice.
\end{description}
\end{frame}





\begin{frame}[fragile]
\frametitle{Ámbitos}

En {\lmr\LaTeX} los comandos globales, o declaraciones globales, son válidos durante todo el ámbito en el que sean utilizados.

\pause

Un ámbito se delimita mediante:

\begin{description}
    \item[Corchetes] \mintinline[escapeinside=||]{latex}{{ <|\textit{ámbito}|> }}
    \item[Grupos] \mintinline[escapeinside=||]{latex}{\begingroup <|\textit{ámbito}|> \endgroup}
    \item[Entornos] \mintinline[escapeinside=||]{latex}{\begin{environment} <|\textit{ámbito}|> \end{environment}}
\end{description}
\end{frame}





\begin{frame}[fragile]
\frametitle{Entornos}

Los entornos son usados para aplicar efectos de composición específicos a una sección del documento.

\pause

Están delimitados mediante las órdenes \mintinline{latex}{\begin{}} y \mintinline{latex}{\end{}}, cuyo argumento obligatorio es el nombre del entorno

\pause

\begin{exampleblock}{}
    \begin{minted}{latex}
\begin{document}
Hello, \LaTeX!
\end{document}
    \end{minted}
\end{exampleblock}
\end{frame}





\section{Tipografía}

\begin{frame}
\frametitle{Familias}
\framesubtitle{Familias, Estilos y Formas}

{\lmr\LaTeX} ofrece una gran variedad de familias tipográficas donde se encuentran:\pause

\begin{itemize}[<+->]
    \item \alert{\bf\lmr Serif}, \alert{\bf\lmr Romana} o \alert{\bf\lmr Con Serifa}.\\
    \hspace*{2ex}Es la fuente por defecto de casi todo documento.
    \item \alert{\bf\lmss Sans Serif}, \alert{\bf\lmss Paloseco} o \alert{\bf\lmss Sin Serifa}.
    \item \alert{\lmtt Typewriter} o \alert{\lmtt mono-espaciada}.
\end{itemize}
\end{frame}


\begin{frame}[fragile]
\frametitle{Familias}
\framesubtitle{Familias, Estilos y Formas}

Para indicar a {\lmr\LaTeX} por cuál de las familias debe cambiar, se pueden utilizar los comandos siguientes:

\begin{center}
    \begin{tabular}{cc}
        \multirow{3}{*}{\lmr Serif}
        & \mintinline{latex}{\textrm{texto}} \\
        & \mintinline{latex}{{\rm texto}} \\
        & \mintinline{latex}{{\rmfamily texto}} \\ \hline
        %
        \multirow{3}{*}{\lmss Sans Serif}
        & \mintinline{latex}{\textsf{texto}} \\
        & \mintinline{latex}{{\sf texto}} \\
        & \mintinline{latex}{{\sffamily texto}} \\ \hline
        %
        \multirow{3}{*}{\lmtt Typewriter}
        & \mintinline{latex}{\texttt{texto}} \\
        & \mintinline{latex}{{\tt texto}} \\
        & \mintinline{latex}{{\ttfamily texto}} \\
    \end{tabular}
\end{center}
\end{frame}




\begin{frame}
\frametitle{Estilos y Formas}
\framesubtitle{Familias, Estilos y Formas}

{\lmr\LaTeX} también ofrece una gran variedad de formas y estilos tipográficos donde se encuentran:

\begin{itemize}[<+->]
    \item \alert{\lmr Upright} o \alert{\lmr Derecha}.
    \item \alert{\bf\lmr Bold} o \alert{\bf\lmr Negrita}.
    \item \alert{\it\lmr Italic} o \alert{\it\lmr Itálica}.
    \item \alert{\sl\lmr Slanted} o \alert{\sl\lmr Inclinada}.
    \item \alert{\sc\lmr Small Caps} o \alert{\sc\lmr Versalitas}.
\end{itemize}
\end{frame}



\begin{frame}[fragile]{Estilos y Formas}{Familias, Estilos y Formas}
Para indicar a {\lmr\LaTeX} por cuál de los estilos o las formas debe cambiar, se pueden utilizar los comandos siguientes:

\begin{center}
    \begin{tabular}{cc}
        \multirow{2}{*}{\lmr Upright}
        & \mintinline{latex}{\textup{texto}} \\
        & \mintinline{latex}{{\upshape texto}} \\ \hline
        %
        \multirow{3}{*}{\bf\lmr Bold}
        & \mintinline{latex}{\textrm{texto}} \\
        & \mintinline{latex}{{\bf texto}} \\
        & \mintinline{latex}{{\bfseries texto}} \\ \hline
        %
        \multirow{3}{*}{\it\lmr Italic}
        & \mintinline{latex}{\textit{texto}} \\
        & \mintinline{latex}{{\it texto}} \\
        & \mintinline{latex}{{\itshape texto}} \\
    \end{tabular}
\end{center}
\end{frame}

\begin{frame}[fragile]{Estilos y Formas}{Familias, Estilos y Formas}
Para indicar a {\lmr\LaTeX} por cuál de los estilos o las formas debe cambiar, se pueden utilizar los comandos siguientes:

\begin{center}
    \begin{tabular}{cc}
        \multirow{3}{*}{\sl\lmr Slanted}
        & \mintinline{latex}{\textsl{texto}} \\
        & \mintinline{latex}{{\sl texto}} \\
        & \mintinline{latex}{{\slshape texto}} \\ \hline
        %
        \multirow{3}{*}{\sc\lmr Small Caps}
        & \mintinline{latex}{\textsc{texto}} \\
        & \mintinline{latex}{{\sc texto}} \\
        & \mintinline{latex}{{\scshape texto}} \\
    \end{tabular}
\end{center}
\end{frame}





\begin{frame}[fragile]
\frametitle{Énfasis}
\framesubtitle{Familias, Estilos y Formas}

Se puede aclarar que las combinaciones de familias, estilos y formas son válidas, y en la mayoría de los casos, conmutativas, pero no todas son posibles.\pause

Si necesitamos enfatizar una parte del texto, lo más apropiado es utilizar los comandos \mintinline{latex}{{\em texto}} o \mintinline{latex}{\emph{texto}}, ya que estas dependen de cada familia, estilo o forma. Por lo que en cada caso {\lmr\LaTeX} escogerá la mejor opción.
\end{frame}

\section{Tamaño de fuente}
\begin{frame}[fragile]
\frametitle{Tamaño de la fuente}

En {\lmr\LaTeX} existen varias formas de cambiar el tamaño de la fuente.\pause

Los más básicos son los siguientes comandos:
\begin{center}
    \begin{tabular}{cc}
        \mintinline{latex}{\tiny} & {\tiny\lmr Muestra} \\
        \mintinline{latex}{\scriptsize} & {\scriptsize\lmr Muestra} \\
        \mintinline{latex}{\footnotesize} & {\footnotesize\lmr Muestra} \\
        \mintinline{latex}{\small} & {\small\lmr Muestra} \\
        \mintinline{latex}{\normalsize} & {\normalsize\lmr Muestra} \\
    \end{tabular}\vline
    \begin{tabular}{cc}
        \mintinline{latex}{\large} & {\large\lmr Muestra} \\
        \mintinline{latex}{\Large} & {\Large\lmr Muestra} \\
        \mintinline{latex}{\LARGE} & {\LARGE\lmr Muestra} \\
        \mintinline{latex}{\huge} & {\huge\lmr Muestra} \\
        \mintinline{latex}{\Huge} & {\Huge\lmr Muestra} \\
    \end{tabular}
\end{center}
\end{frame}

\begin{frame}[fragile]
\frametitle{Tamaño de la fuente}
Otra forma para cambiar el tamaño de la fuente del documento es mediante las opciones de \mintinline{latex}{\documentclass}.\pause

Estas opciones pueden ser \texttt{10pt}, \texttt{11pt} o \texttt{12pt}, con el significado habitual de todo editor. Estas opciones son aplicadas de manera general al documento.
\end{frame}



\section{{\lmr\LaTeX} en español}

\begin{frame}[fragile]
\frametitle{{\lmr\LaTeX} en español}
Al escribir en un documento, colocar títulos y demás rótulos, estos tienen el formato e idioma inglés.\pause

Para cambiarlo, y utilizar rótulos en español es necesario incluir el paquete \texttt{babel}.\pause

\begin{minted}{latex}
\usepackage[spanish]{babel}
\end{minted}

\pause

Para seguir lo dicho por la RAE ASALE en Ortografía de la Lengua Española (2010), utilizaremos la opción \texttt{es-nodecimaldot}.


\begin{minted}{latex}
\usepackage[spanish,es-nodecimaldot]{babel}
\end{minted}
\end{frame}





\section{Listas}

\subsection{Listas Desordenadas}


\begin{frame}[fragile]
\frametitle{Listas Desordenadas}

Para crear listas desordenadas, o listas de viñetas, utilizamos el entorno \texttt{itemize}.

\pause

Utilizaremos el comando \mintinline{latex}{\item} para generar nuevos elementos en la lista.

\pause

\begin{columns}
    \begin{column}{0.45\textwidth}
        \begin{block}{}
            \begin{minted}{latex}
\begin{itemize}
    \item Pan
    \item Jamón
    \item Queso
\end{itemize}
            \end{minted}
        \end{block}
    \end{column}
    \begin{column}{0.1\textwidth}
        \centering $\rightarrow$
    \end{column}
    \begin{column}{0.45\textwidth}
        \begin{block}{}
            \lmr
            \begin{tabular}{cl}
                $\bullet$ & Pan \\
                $\bullet$ & Jamón \\
                $\bullet$ & Queso \\
            \end{tabular}
        \end{block}
    \end{column}
\end{columns}
\end{frame}





\subsection{Listas Ordenadas}

\begin{frame}[fragile]
\frametitle{Listas Ordenadas}

Para crear listas ordenadas, o listas numeradas, utilizamos el entorno \texttt{enumarate}.

\pause

Utilizaremos el comando \mintinline{latex}{\item} para generar nuevos elementos en la lista.

\pause

\begin{columns}
    \begin{column}{0.45\textwidth}
        \begin{block}{}
            \begin{minted}{latex}
\begin{enumerate}
    \item Colocar pan
    \item Añadir jamón
    \item Rayar queso
\end{enumerate}
            \end{minted}
        \end{block}
    \end{column}
    \begin{column}{0.1\textwidth}
        \centering $\rightarrow$
    \end{column}
    \begin{column}{0.45\textwidth}
        \begin{block}{}
            \lmr
            \begin{tabular}{cl}
                1. & Colocar pan \\
                2. & Añadir jamón \\
                3. & Rayar queso \\
            \end{tabular}
        \end{block}
    \end{column}
\end{columns}
\end{frame}





\subsection{Anidación}

\begin{frame}[fragile]
\frametitle{Anidación}

Anidar listas es bastante sencillo. Solamente debemos colocar un entorno dentro del otro.

\pause

\begin{columns}
    \begin{column}{0.45\textwidth}
        \begin{block}{}
            \begin{minted}{latex}
\begin{itemize}
    \item Pan
    \item Jamón
    \begin{enumerate}
        \item Pietran
        \item Vienna
    \end{enumerate}
    \item Queso
\end{itemize}
            \end{minted}
        \end{block}
    \end{column}
    \begin{column}{0.1\textwidth}
        \centering $\rightarrow$
    \end{column}
    \begin{column}{0.45\textwidth}
        \begin{block}{}
            \lmr
            \begin{tabular}{cl}
                $\bullet$ & Pan \\
                $\bullet$ & Jamón \\
                & \begin{tabular}{cl}
                    1. & Pietran \\
                    2. & Vienna \\
                \end{tabular} \\
                $\bullet$ & Queso \\
            \end{tabular}
        \end{block}
    \end{column}
\end{columns}
\end{frame}





\section{Tablas}

\begin{frame}[fragile]
\frametitle{Tablas}

{\lmr\LaTeX} tiene un entorno muy cómodo para la creación de tablas. \texttt{tabular}.

\pause

\begin{columns}
    \begin{column}{0.45\textwidth}
        \begin{block}{}
            \begin{minted}[escapeinside=||]{latex}
\begin{tabular}
    [<|\textit{posición}|>]
    {<|\textit{formato}|>}
    ... & ... & ... \\
    ... & ... & ... \\
    ... & ... & ... \\
\end{tabular}
            \end{minted}
        \end{block}
    \end{column}
    \pause
    \begin{column}{0.1\textwidth}
        \centering $\rightarrow$
    \end{column}
    \begin{column}{0.45\textwidth}
        \begin{block}{}
            \lmr
            \begin{tabular}{ccc}
                \ldots & \ldots & \ldots \\
                \ldots & \ldots & \ldots \\
                \ldots & \ldots & \ldots \\
            \end{tabular}
        \end{block}
    \end{column}
\end{columns}
\end{frame}





\begin{frame}[fragile]
\frametitle{Tablas}

Existe una versión estrella de este entorno que nos permite definir tablas con un ancho fijo predeterminado. \texttt{tabular*}.

\pause

\begin{block}{}
    \begin{minted}[escapeinside=||]{latex}
\begin{tabular*}
    {<|\textit{ancho}|>}
    {<|\textit{formato}|>}
    ... & ... & ... \\
    ... & ... & ... \\
    ... & ... & ... \\
\end{tabular*}
    \end{minted}
\end{block}
\end{frame}





\begin{frame}
\frametitle{Tablas}

El argumento \texttt{<formato>} contiene la información sobre el número de columnas y su justificación.

\begin{description}
    \item[\texttt{l}] left --- Izquierda
    \item[\texttt{c}] center --- Centrada
    \item[\texttt{r}] right --- Derecha
    \item[\texttt{p{<ancho>}}] paragraph --- Columna de ancho fijo predefinido con párrafos
    \item[\texttt{|}] Línea Vertical
\end{description}
\end{frame}





\begin{frame}[fragile]
\frametitle{Tablas}
\begin{columns}
    \begin{column}{0.45\textwidth}
        \begin{block}{}
            \begin{minted}{latex}
\begin{tabular}{lcr|p{2.3cm}}
    a & b & c & lorem... \\
    aa & bb & cc & lorem... \\
    aaa & bbb & ccc & lorem.. \\
\end{tabular}
            \end{minted}
        \end{block}
    \end{column}
    \begin{column}{0.1\textwidth}
        \centering $\rightarrow$
    \end{column}
    \begin{column}{0.45\textwidth}
        \begin{block}{}
            \lmr
            \begin{tabular}{lcr|p{2.3cm}}
                a & b & c & lorem ipsum dolor sit ame \\
                aa & bb & cc & lorem ipsum dolor sit ame \\
                aaa & bbb & ccc & lorem ipsum dolo sit amet \\
            \end{tabular}
        \end{block}
    \end{column}
\end{columns}
\end{frame}





\begin{frame}[fragile]
\frametitle{Tablas}
\begin{columns}
    \begin{column}{0.45\textwidth}
        \begin{block}{}
            \begin{minted}{latex}
\begin{tabular}{lll}
    \hline
    a & b & c \\
    \hline
    aa & bb & cc \\
    aaa & bbb & ccc \\
    \hline
\end{tabular}
            \end{minted}
        \end{block}
    \end{column}
    \begin{column}{0.1\textwidth}
        \centering $\rightarrow$
    \end{column}
    \begin{column}{0.45\textwidth}
        \begin{block}{}
            \lmr
            \begin{tabular}{lll}
                \hline
                a & b & c \\
                \hline
                aa & bb & cc \\
                aaa & bbb & ccc \\
                \hline
            \end{tabular}
        \end{block}
    \end{column}
\end{columns}
\end{frame}





\begin{frame}[fragile]
\frametitle{Tablas}
\begin{columns}
    \begin{column}{0.45\textwidth}
        \begin{block}{}
            \begin{minted}{latex}
\begin{tabular}{|l||ll|}
    \hline
    a & b & c \\
    \hline
    \hline
    aa & bb & cc \\
    aaa & bbb & ccc \\
    \hline
\end{tabular}
            \end{minted}
        \end{block}
    \end{column}
    \begin{column}{0.1\textwidth}
        \centering $\rightarrow$
    \end{column}
    \begin{column}{0.45\textwidth}
        \begin{block}{}
            \lmr
            \begin{tabular}{|l||ll|}
                \hline
                a & b & c \\
                \hline
                \hline
                aa & bb & cc \\
                aaa & bbb & ccc \\
                \hline
            \end{tabular}
        \end{block}
    \end{column}
\end{columns}
\end{frame}





\begin{frame}[fragile]
\frametitle{Tablas}

Todas las tablas en {\lmr\LaTeX} se comportan como \alert{cajas}

\pause

\begin{columns}
    \begin{column}{0.45\textwidth}
        \begin{block}{}
            \begin{minted}{latex}
texto
\begin{tabular}{|c|}
    \hline
    a \\
    aa \\
    \hline
\end{tabular}
texto
            \end{minted}
        \end{block}
    \end{column}
    \begin{column}{0.1\textwidth}
        \centering $\rightarrow$
    \end{column}
    \begin{column}{0.45\textwidth}
        \begin{block}{}
            \lmr
            texto
            \begin{tabular}{|c|}
                \hline
                a \\
                aa \\
                \hline
            \end{tabular}
            texto
        \end{block}
    \end{column}
\end{columns}
\end{frame}





\begin{frame}[fragile]
\frametitle{Tablas}

Todas las tablas en {\lmr\LaTeX} se comportan como \alert{cajas}

\begin{columns}
    \begin{column}{0.45\textwidth}
        \begin{block}{}
            \begin{minted}{latex}
texto
\begin{tabular}[t]{|c|}
    \hline
    a \\
    aa \\
    \hline
\end{tabular}
texto
            \end{minted}
        \end{block}
    \end{column}
    \begin{column}{0.1\textwidth}
        \centering $\rightarrow$
    \end{column}
    \begin{column}{0.45\textwidth}
        \begin{block}{}
            \lmr
            texto
            \begin{tabular}[t]{|c|}
                \hline
                a \\
                aa \\
                \hline
            \end{tabular}
            texto
        \end{block}
    \end{column}
\end{columns}
\end{frame}





\begin{frame}[fragile]
\frametitle{Tablas}

Todas las tablas en {\lmr\LaTeX} se comportan como \alert{cajas}

\begin{columns}
    \begin{column}{0.45\textwidth}
        \begin{block}{}
            \begin{minted}{latex}
texto
\begin{tabular}[b]{|c|}
    \hline
    a \\
    aa \\
    \hline
\end{tabular}
texto
            \end{minted}
        \end{block}
    \end{column}
    \begin{column}{0.1\textwidth}
        \centering $\rightarrow$
    \end{column}
    \begin{column}{0.45\textwidth}
        \begin{block}{}
            \lmr
            texto
            \begin{tabular}[b]{|c|}
                \hline
                a \\
                aa \\
                \hline
            \end{tabular}
            texto
        \end{block}
    \end{column}
\end{columns}
\end{frame}





\subsection{Celdas Complejas}

\begin{frame}[fragile]
\frametitle{Celdas Complejas}
\framesubtitle{\texttt{multicolumn}}

Para combinar celdas horizontalmente usamos el comando \mintinline{latex}{multicolumn}:

\begin{minted}[escapeinside=||]{latex}
\multicolumn{<|\textit{columnas}|>}{<|\textit{formato}|>}{<|\textit{texto}|>}
\end{minted}
\end{frame}



\begin{frame}[fragile]
\frametitle{Celdas Complejas}
\framesubtitle{\texttt{multicolumn}}

\begin{onlyenv}<1>
    \begin{minted}{latex}
\begin{tabular}{lrr}
    \hline
    \multicolumn{3}{c}{Número de animales} \\
    \hline
    & Macho & Hembra \\
    León & 1 & 12 \\
    Tigre & 4 & 10 \\
    Pingüino & 15 & 15 \\
    \hline
\end{tabular}
    \end{minted}
\end{onlyenv}


\begin{onlyenv}<2->
    \lmr
    \begin{tabular}{lrr}
        \hline
        \multicolumn{3}{c}{Número de animales} \\
        \hline
        & Macho & Hembra \\
        León & 1 & 12 \\
        Tigre & 4 & 10 \\
        Pingüino & 15 & 15 \\
        \hline
    \end{tabular}
\end{onlyenv}
\end{frame}





\begin{frame}[fragile]
\frametitle{Celdas Complejas}
\framesubtitle{\texttt{multicolumn}}

Para combinar celdas verticalmente usamos el comando \mintinline{latex}{multirow} del paquete hómonimo:

\begin{minted}[escapeinside=||]{latex}
\multirow[<|\textit{vpos}|>]{<|\textit{filas}|>}{<|\textit{ancho}|>}{<|\textit{text}|>}
\end{minted}

\pause

\begin{minted}{latex}
\usepackage{multirow}
\end{minted}
\end{frame}





\begin{frame}[fragile]
\frametitle{Celdas Complejas}
\framesubtitle{\texttt{multicolumn}}

\begin{onlyenv}<1>
    \begin{minted}{latex}
\begin{tabular}{clrr}
    \hline
    \multirow{4}{*}{Número de animales}
    && Macho & Hembra \\
    & Leon & 1 & 12 \\
    & Tigre & 4 & 10 \\
    & Pingüino & 15 & 15 \\
    \hline
\end{tabular}
    \end{minted}
\end{onlyenv}
\begin{onlyenv}<2->
    \lmr
        \begin{tabular}{clrr}
            \hline
            \multirow{4}{*}{Número de animales}
            && Macho & Hembra \\
            & Leon & 1 & 12 \\
            & Tigre & 4 & 10 \\
            & Pingüino & 15 & 15 \\
            \hline
        \end{tabular}
\end{onlyenv}
\end{frame}





\begin{frame}
\frametitle{Tablas}

\begin{alertblock}{IMPORTANTE}
    El entorno \texttt{tabular} hace que {\lmr\LaTeX} entre en el modo \alert{ID}
\end{alertblock}
\end{frame}





\section{Cajas}

\begin{frame}
\frametitle{Cajas}

Una de las herramientas de edición más básicas que existen en {\lmr\LaTeX} son las \alert{cajas}.

\pause

{\lmr\LaTeX} trata las cajas siempre como una sola unidad, como una letra. Por lo que su contenido nunca será dividido.

\pause

\begin{alertblock}{IMPORTANTE}
    Todas las cajas hacen que {\lmr\LaTeX} entre en el modo \alert{ID}
\end{alertblock}
\end{frame}





\begin{frame}[fragile]
\frametitle{Cajas}

\begin{columns}
    \begin{column}{0.45\textwidth}
        \begin{block}{}
            \begin{minted}{latex}
\mbox{Texto}

\fbox{Texto}
            \end{minted}
        \end{block}
    \end{column}
    \begin{column}{0.1\textwidth}
        \centering $\rightarrow$
    \end{column}
    \begin{column}{0.45\textwidth}
        \begin{block}{}
            \lmr
            \mbox{Texto}

            \bigskip

            \fbox{Texto}
        \end{block}
    \end{column}
\end{columns}
\end{frame}





\begin{frame}[fragile]
\frametitle{Cajas}

\begin{columns}
    \begin{column}{0.45\textwidth}
        \begin{block}{}
            \begin{minted}{latex}
\makebox[2cm]{Texto}

\framebox[2cm]{Texto}
            \end{minted}
        \end{block}
    \end{column}
    \begin{column}{0.1\textwidth}
        \centering $\rightarrow$
    \end{column}
    \begin{column}{0.45\textwidth}
        \begin{block}{}
            \lmr
            \makebox[2cm]{Texto}

            \bigskip

            \framebox[2cm]{Texto}
        \end{block}
    \end{column}
\end{columns}
\end{frame}





\begin{frame}[fragile]
\frametitle{Cajas}

\begin{columns}
    \begin{column}{0.45\textwidth}
        \begin{block}{}
            \begin{minted}{latex}
\makebox[2cm][l]{Texto}

\framebox[2cm][l]{Texto}
            \end{minted}
        \end{block}
    \end{column}
    \begin{column}{0.1\textwidth}
        \centering $\rightarrow$
    \end{column}
    \begin{column}{0.45\textwidth}
        \begin{block}{}
            \lmr
            \makebox[2cm][l]{Texto}

            \bigskip

            \framebox[2cm][l]{Texto}
        \end{block}
    \end{column}
\end{columns}
\end{frame}





\begin{frame}[fragile]
\frametitle{Cajas}

\begin{columns}
    \begin{column}{0.45\textwidth}
        \begin{block}{}
            \begin{minted}{latex}
\makebox[2cm][r]{Texto}

\framebox[2cm][r]{Texto}
            \end{minted}
        \end{block}
    \end{column}
    \begin{column}{0.1\textwidth}
        \centering $\rightarrow$
    \end{column}
    \begin{column}{0.45\textwidth}
        \begin{block}{}
            \lmr
            \makebox[2cm][r]{Texto}

            \bigskip

            \framebox[2cm][r]{Texto}
        \end{block}
    \end{column}
\end{columns}
\end{frame}





\begin{frame}[fragile]
\frametitle{Cajas}

\begin{columns}
    \begin{column}{0.45\textwidth}
        \begin{block}{}
            \begin{minted}{latex}
\makebox[4cm][s]{Texto texto}

\framebox[4cm][s]{Texto texto}
            \end{minted}
        \end{block}
    \end{column}
    \begin{column}{0.1\textwidth}
        \centering $\rightarrow$
    \end{column}
    \begin{column}{0.45\textwidth}
        \begin{block}{}
            \lmr
            \hbadness=10000
            \makebox[4cm][s]{Texto texto}

            \bigskip

            \framebox[4cm][s]{Texto texto}
        \end{block}
    \end{column}
\end{columns}
\end{frame}





\begin{frame}
\frametitle{Dimensiones}

Podemos usar cualquier dimensión de {\lmr\TeX} para definir medidas.

\begin{center}
\scriptsize
\begin{tabular}{>{\ttfamily}lll}
    \toprule
    Símbolo & Unidad & Definición \\
    \midrule
    in & Pulgada & \\
    cm & Centímetro & 2.54 cm = 1 in \\
    mm & Milimetro & 10 mm = 1 cm \\
    pt & Punto & 72.27 pt = 1 in \\
    pc & Pica & 1 pc = 12 pt \\
    bp & Punto Grande & 72 bp = 1 in \\
    dd & Punto Didot & 1157 dd = 1238 pt \\
    cc & Cicero & 1 cc = 12 dd \\
    sp & Punto Escalado & 65536 sp = 1 pt \\
    ex & Equis & El alto de la \texttt{x} en la fuente actual \\
    em & Eme & El ancho de la \texttt{M} en la fuente actual \\
    fil & Infinito de Primer Orden & \(\aleph_0\) \\
    fill & Infinito de Segundo Orden & \(\aleph_1\) \\
    filll & Infinito de Tercer Orden & \(\aleph_2\) \\
    \bottomrule
\end{tabular}
\end{center}
\end{frame}





\begin{frame}[fragile]
\frametitle{Dimensiones}

\begin{columns}
    \begin{column}{0.45\textwidth}
        \begin{block}{}
            \begin{minted}{latex}
\framebox[20ex][c]{Texto texto}
            \end{minted}
        \end{block}
    \end{column}
    \begin{column}{0.1\textwidth}
        \centering $\rightarrow$
    \end{column}
    \begin{column}{0.45\textwidth}
        \begin{block}{}
            \lmr
            \begin{tikzpicture}
                \node [inner sep=0, anchor=west] {\framebox[20ex][c]{Texto texto}};

                \draw (0ex, 1.2em) node [inner sep=0, anchor=south, rotate=-90] {x}
                    ++(1ex, 0.3ex) node [inner sep=0, anchor=south, rotate=-90] {x}
                    ++(1ex,-0.3ex) node [inner sep=0, anchor=south, rotate=-90] {x}
                    ++(1ex, 0.3ex) node [inner sep=0, anchor=south, rotate=-90] {x}
                    ++(1ex,-0.3ex) node [inner sep=0, anchor=south, rotate=-90] {x}
                    ++(1ex, 0.3ex) node [inner sep=0, anchor=south, rotate=-90] {x}
                    ++(1ex,-0.3ex) node [inner sep=0, anchor=south, rotate=-90] {x}
                    ++(1ex, 0.3ex) node [inner sep=0, anchor=south, rotate=-90] {x}
                    ++(1ex,-0.3ex) node [inner sep=0, anchor=south, rotate=-90] {x}
                    ++(1ex, 0.3ex) node [inner sep=0, anchor=south, rotate=-90] {x}
                    ++(1ex,-0.3ex) node [inner sep=0, anchor=south, rotate=-90] {x}
                    ++(1ex, 0.3ex) node [inner sep=0, anchor=south, rotate=-90] {x}
                    ++(1ex,-0.3ex) node [inner sep=0, anchor=south, rotate=-90] {x}
                    ++(1ex, 0.3ex) node [inner sep=0, anchor=south, rotate=-90] {x}
                    ++(1ex,-0.3ex) node [inner sep=0, anchor=south, rotate=-90] {x}
                    ++(1ex, 0.3ex) node [inner sep=0, anchor=south, rotate=-90] {x}
                    ++(1ex,-0.3ex) node [inner sep=0, anchor=south, rotate=-90] {x}
                    ++(1ex, 0.3ex) node [inner sep=0, anchor=south, rotate=-90] {x}
                    ++(1ex,-0.3ex) node [inner sep=0, anchor=south, rotate=-90] {x}
                    ++(1ex, 0.3ex) node [inner sep=0, anchor=south, rotate=-90] {x};
            \end{tikzpicture}
        \end{block}
    \end{column}
\end{columns}
\end{frame}





\section{\texttt{booktabs}}

\begin{frame}[fragile]
\frametitle{Tablas Profesionales con \texttt{booktabs}}

El paquete \texttt{booktabs} mejora significativamente la estética de las tablas.

\pause

Reemplaza \mintinline{latex}{\hline} con comandos con grosor específico

\begin{description}
    \item[\mintinline{latex}{\toprule}] Línea superior gruesa
    \item[\mintinline{latex}{\midrule}] Línea intermedia delgada
    \item[\mintinline{latex}{\bottomrule}] Línea inferior gruesa
\end{description}
\end{frame}





\begin{frame}[fragile]
\frametitle{Tablas Profesionales con \texttt{booktabs}}

\begin{columns}
    \begin{column}{0.45\textwidth}
        \begin{block}{}
            \begin{minted}{latex}
\begin{tabular}{llc}
    \toprule
    Partícula &
    Descubridor &
    Año \\
    \midrule
    Electrón & Thompson & 1987 \\
    Protón & Rutherford & 1919 \\
    Neutrón & Chadwick & 1932 \\
    \bottomrule
\end{tabular}
            \end{minted}
        \end{block}
    \end{column}
    \begin{column}{0.1\textwidth}
        \centering $\rightarrow$
    \end{column}
    \begin{column}{0.45\textwidth}
        \begin{block}{}
            \lmr
            \begin{tabular}{llc}
                \toprule
                    Partícula & Descubridor & Año \\
                    \midrule
                    Electrón & Thompson & 1987 \\
                    Protón & Rutherford & 1919 \\
                    Neutrón & Chadwick & 1932 \\
                \bottomrule
            \end{tabular}
        \end{block}
    \end{column}
\end{columns}
\end{frame}





\section{\texttt{siunitx}}

\begin{frame}[fragile]
\frametitle{Manejo de Unidades y Números con \texttt{siunitx}}

El paquete \texttt{siunitx} garantiza un formato consistente para números y unidades del Sistema Internacional (SI).

\pause

\begin{columns}
    \begin{column}{0.45\textwidth}
        \begin{block}{}
            \begin{minted}{latex}
\num{12523754.32}

\unit{\metre\per\second\squared}

\qty{735.43}{\mega\electronvolt}

\qty{28.7 +- 0.5}{\kilo\newton}
            \end{minted}
        \end{block}
    \end{column}
    \begin{column}{0.1\textwidth}
        \centering $\rightarrow$
    \end{column}
    \begin{column}{0.45\textwidth}
        \begin{block}{}
            \lmr
            \num{12523754.32}

            \bigskip

            \unit{\metre\per\second\squared}

            \bigskip

            \qty{735.43}{\mega\electronvolt}

            \bigskip

            \qty{28.7 +- 0.5}{\kilo\newton}
        \end{block}
    \end{column}
\end{columns}
\end{frame}





\begin{frame}[fragile]
\frametitle{Manejo de Unidades y Números con \texttt{siunitx}}

El paquete \texttt{siunitx} garantiza un formato consistente para números y unidades del Sistema Internacional (SI).

\begin{columns}
    \begin{column}{0.45\textwidth}
        \begin{block}{}
            \begin{minted}{latex}
\num{12523754.32}

\unit{\metre\per\second\squared}

\qty{735.43}{\mega\electronvolt}

\sisetup{uncertainty-mode=separate}

\qty{28.7 +- 0.5}{\kilo\newton}
            \end{minted}
        \end{block}
    \end{column}
    \begin{column}{0.1\textwidth}
        \centering $\rightarrow$
    \end{column}
    \begin{column}{0.45\textwidth}
        \begin{block}{}
            \lmr
            \num{12523754.32}

            \bigskip

            \unit{\metre\per\second\squared}

            \bigskip

            \qty{735.43}{\mega\electronvolt}

            \bigskip
            \sisetup{uncertainty-mode=separate}

            \qty{28.7 +- 0.5}{\kilo\newton}
        \end{block}
    \end{column}
\end{columns}
\end{frame}




\begin{frame}[fragile]
\frametitle{Manejo de Unidades y Números con \texttt{siunitx}}

\begin{columns}
    \begin{column}{0.45\textwidth}
        \begin{block}{}
            \begin{minted}{latex}
\begin{tabular}
    {S[table-format=2.3]}
    12.345 \\
    0.1 \\
    99.99 \\
\end{tabular}
            \end{minted}
        \end{block}
    \end{column}
    \begin{column}{0.1\textwidth}
        \centering $\rightarrow$
    \end{column}
    \begin{column}{0.45\textwidth}
        \begin{block}{}
            \lmr
            \begin{tabular}{S[table-format=2.3]}
                12.345 \\
                0.1 \\
                99.99 \\
            \end{tabular}
        \end{block}
    \end{column}
\end{columns}
\end{frame}





\begin{frame}[fragile]
\frametitle{Manejo de Unidades y Números con \texttt{siunitx}}

\begin{columns}
    \begin{column}{0.45\textwidth}
        \begin{block}{}
            \begin{minted}{latex}
\begin{tabular}
    {S[table-format=2.3(1)]}
    12.345 +- 0.001 \\
    0.1 +- 0.1\\
    99.99 +- 0.01 \\
\end{tabular}
            \end{minted}
        \end{block}
    \end{column}
    \begin{column}{0.1\textwidth}
        \centering $\rightarrow$
    \end{column}
    \begin{column}{0.45\textwidth}
        \begin{block}{}
            \lmr
            \begin{tabular}{S[table-format=2.3(1)]}
                12.345 +- 0.001 \\
                0.1 +- 0.1\\
                99.99 +- 0.01 \\
            \end{tabular}
        \end{block}
    \end{column}
\end{columns}
\end{frame}





\begin{frame}[fragile]
\frametitle{Manejo de Unidades y Números con \texttt{siunitx}}

\begin{columns}
    \begin{column}{0.45\textwidth}
        \begin{block}{}
            \begin{minted}{latex}
\begin{tabular}
    {S[table-format=2.3(1),
    uncertainty-mode=separate]}
    12.345 +- 0.001\\
    0.1 +- 0.1 \\
    99.99 +- 0.01 \\
\end{tabular}
            \end{minted}
        \end{block}
    \end{column}
    \begin{column}{0.1\textwidth}
        \centering $\rightarrow$
    \end{column}
    \begin{column}{0.45\textwidth}
        \begin{block}{}
            \lmr
            \begin{tabular}{S[table-format=2.3(1),uncertainty-mode=separate]}
                12.345 +- 0.001\\
                0.1 +- 0.1 \\
                99.99 +- 0.01 \\
            \end{tabular}
        \end{block}
    \end{column}
\end{columns}
\end{frame}





\begin{frame}[allowframebreaks]
\frametitle{Referencias}
\nocite{*}
\printbibliography
\end{frame}
\end{document}
